% bie_center.tex
%
% Purpose:
%   Present the center-obstacle boundary integral representation and the full
%   port trace-matching matrix used for lead-scattering and eigenvalue problems.
%
% Main algorithm:
%   Combine the Muller BIE row on Sigma with two Cauchy trace matching rows at
%   Gamma_- and Gamma_+, using center homogeneous Rayleigh coefficients and
%   outgoing half-lead trace bases.
%
% Based on:
%   notes/theory/scat_formulation2.md.
%
% Main changes:
%   The derivation is rewritten in English and uses the newer center-cell
%   notation a,b instead of c^+,c^-.
%
% Numerical goal:
%   Record the block matrix \mathcal A and the dimension count that should be
%   checked in implementations.

\section{Center-Cell Boundary Integral and Port Matching}

\subsection{Center homogeneous field}

In the center background region, the homogeneous Rayleigh part is written
\[
  u_h(x,y)
  =
  \sum_{m=-M}^{M}
  a_m\ee^{\ii\gamma_m(x-X_-)}\psi_m(y)
  +
  \sum_{m=-M}^{M}
  b_m\ee^{-\ii\gamma_m(x-X_-)}\psi_m(y).
\]
The vectors \(a,b\in\CC^K\) are local right-going and left-going Rayleigh
coefficients in the center cell.  They are not positive and negative lead
indices.

Let
\[
  E=\diag(\ee^{\ii\gamma_mL_0}),
  \qquad
  L_0=X_+-X_-,
  \qquad
  \GammaR=\diag(\gamma_m).
\]
With the center-cell outward normals, the port traces of \(u_h\) are
\[
  D_h^- = a+b,
  \qquad
  N_h^- = -\ii\GammaR(a-b),
\]
and
\[
  D_h^+ = Ea+E^{-1}b,
  \qquad
  N_h^+ = \ii\GammaR(Ea-E^{-1}b).
\]

\subsection{Muller representation on the center obstacle}

The exterior center-cell field is decomposed as
\[
  u_e=u_h+u_s^e,
  \qquad
  x\in\Ccell^0\setminus\overline{\Omega_0},
\]
while the interior field is represented by
\[
  u_i=u_s^i,
  \qquad
  x\in\Omega_0 .
\]
The density convention used by the current code is
\[
  \eta=
  \begin{bmatrix}
  \tau\\
  -\sigma
  \end{bmatrix}.
\]
The scattered fields are represented by layer potentials such as
\[
  u_s^e
  =
  D_{\mathrm{QP}}^{(k_e)}[\tau]
  +
  S_{\mathrm{QP}}^{(k_e)}[\sigma],
\]
with the corresponding interior representation using \(k_i\).  The assembled
Muller matrix on \(\Sigma=\partial\Omega_0\) is denoted \(\Aqp\).

On \(\Sigma\), the transmission conditions for \(u_e=u_h+u_s^e\) and \(u_i=u_s^i\)
give the forced Muller row
\[
  \Aqp\eta+H_\Sigma^a a+H_\Sigma^b b=0 .
\]
Here \(H_\Sigma^a\) and \(H_\Sigma^b\) are the Cauchy data on \(\Sigma\) of the
center homogeneous Rayleigh basis.  For boundary nodes
\[
  z_j=(x_j,y_j)\in\Sigma,
  \qquad
  \nu_j=(\nu_{x,j},\nu_{y,j}),
\]
their Dirichlet blocks are
\[
  H_{D,\Sigma}^{a}(j,m)
  =
  \ee^{\ii\gamma_m(x_j-X_-)}\psi_m(y_j),
  \qquad
  H_{D,\Sigma}^{b}(j,m)
  =
  \ee^{-\ii\gamma_m(x_j-X_-)}\psi_m(y_j),
\]
and their Neumann blocks are
\[
  H_{N,\Sigma}^{a}(j,m)
  =
  \ii(\gamma_m\nu_{x,j}+\beta_m\nu_{y,j})
  \ee^{\ii\gamma_m(x_j-X_-)}\psi_m(y_j),
\]
\[
  H_{N,\Sigma}^{b}(j,m)
  =
  \ii(-\gamma_m\nu_{x,j}+\beta_m\nu_{y,j})
  \ee^{-\ii\gamma_m(x_j-X_-)}\psi_m(y_j).
\]

The center homogeneous coefficients \(a,b\) are essential.  The density
\(\eta\) represents the scattered correction generated by the center inclusion;
it does not span the full homogeneous solution space in the center cell.  If
\(a,b\) are omitted, the port matching system is generally overdetermined and
does not represent the intended center-cell solution class.

\subsection{Density-generated outgoing Rayleigh coefficients}

The center density produces outgoing Rayleigh coefficients at the two ports:
\[
  s^- = F_-\eta,
  \qquad
  s^+ = F_+\eta .
\]
These maps are implemented by \texttt{bloch.farfield\_extractors}.  They are
linear maps from the boundary density to outgoing Rayleigh coefficients, not new
unknown vectors.  With the present outward-normal convention, the density
contribution to the port traces is
\[
  D_s^- = F_-\eta,
  \qquad
  N_s^- = \ii\GammaR F_-\eta,
\]
and
\[
  D_s^+ = F_+\eta,
  \qquad
  N_s^+ = \ii\GammaR F_+\eta.
\]
The same sign \(+\ii\GammaR\) appears on both sides because the left outgoing
Rayleigh factor has \(\partial_x=-\ii\gamma_m\), while the left port outward
normal is \(-e_x\).

\subsection{Port matching equations}

The center trace is the sum of the homogeneous part and the density-generated
scattered part.  The half-lead trace is the sum of prescribed incoming data and
unknown outgoing data:
\[
  D_h^\pm+D_s^\pm
  =
  \Din^\pm p_{\mathrm{in}}^\pm+\Dout^\pm r^\pm,
\]
\[
  N_h^\pm+N_s^\pm
  =
  \Nin^\pm p_{\mathrm{in}}^\pm+\Nout^\pm r^\pm .
\]
After substituting the formulas above, the full unknown vector is
\[
  x=
  \begin{bmatrix}
  \eta\\
  a\\
  b\\
  r^-\\
  r^+
  \end{bmatrix}.
\]
The coupled matrix is
\[
  \Atr
  =
  \begin{bmatrix}
  \Aqp & H_\Sigma^a & H_\Sigma^b & 0 & 0\\
  F_- & I & I & -\Dout^- & 0\\
  \ii\GammaR F_- & -\ii\GammaR & \ii\GammaR & -\Nout^- & 0\\
  F_+ & E & E^{-1} & 0 & -\Dout^+\\
  \ii\GammaR F_+ & \ii\GammaR E & -\ii\GammaR E^{-1} & 0 & -\Nout^+
  \end{bmatrix}.
\]
For forced lead scattering,
\[
  \Atr x=b_{\mathrm{in}},
\]
where
\[
  b_{\mathrm{in}}
  =
  \begin{bmatrix}
  0\\
  \Din^-p_{\mathrm{in}}^-\\
  \Nin^-p_{\mathrm{in}}^-\\
  \Din^+p_{\mathrm{in}}^+\\
  \Nin^+p_{\mathrm{in}}^+
  \end{bmatrix}.
\]
For the eigenvalue or TEP-type workflow, one uses
\[
  \Atr x=0 .
\]

\subsection{Dimension count}

Let \(N=N_{\mathrm{tot}}\) denote the number of boundary quadrature nodes on
\(\Sigma\), so that \(\eta\in\CC^{2N}\).  The BIE/transmission row contributes
\(2N\) equations.  Each port contributes \(K\) Dirichlet and \(K\) Neumann trace
matching equations, so the two ports contribute \(4K\) equations.  The total
number of equations is therefore
\[
  2N+4K .
\]

The unknowns are
\[
  \eta\in\CC^{2N},
  \qquad
  a,b\in\CC^K,
  \qquad
  r^\pm\in\CC^{K_{\mathrm{out}}^\pm}.
\]
The total number of unknowns is
\[
  2N+2K+K_{\mathrm{out}}^-+K_{\mathrm{out}}^+ .
\]
Thus the matrix is square when
\[
  K_{\mathrm{out}}^-+K_{\mathrm{out}}^+=2K .
\]
In a non-degenerate lossless second-order setting, one expects each outgoing
half-lead trace subspace to be approximately half-dimensional:
\[
  K_{\mathrm{out}}^\pm\approx K .
\]
This is a diagnostic, not an unconditional theorem.  Deviations may indicate
unit-circle propagating modes, band edges, grazing Rayleigh orders, mode
collisions, thresholding issues, or insufficient Rayleigh truncation.
